\chapter{Architecture}
\label{chap-architecture}

\section{A single rule: dependencies point inwards}

The architecture fits in one sentence: \textbf{the game engine does not know how it is
    displayed; the interfaces know the engine; nothing knows the interfaces}
(\cref{fig-couches}).

\begin{figure}[ht]
    \centering
    \begin{tikzpicture}[node distance=7mm and 9mm]
        \node[coeur, minimum width=11.2cm, minimum height=1.5cm] (core)
        {\textbf{battleship\_core} --- game engine (the \emph{model})\\[2pt]
            \footnotesize C11, standard library only, no dynamic allocation};

        \node[couche, above=14mm of core.north west, anchor=south west, minimum width=3.4cm,
            minimum height=1.6cm] (gui)
        {\textbf{battleship}\\[1pt]\footnotesize GTK~3 application\\\footnotesize view + controller};
        \node[couche, above=14mm of core.north, anchor=south, minimum width=3.4cm,
            minimum height=1.6cm] (cli)
        {\textbf{battleship-cli}\\[1pt]\footnotesize terminal game\\\footnotesize simulator};
        \node[couche, above=14mm of core.north east, anchor=south east, minimum width=3.4cm,
            minimum height=1.6cm] (tests)
        {\textbf{tests}\\[1pt]\footnotesize unit\\\footnotesize and property};

        \node[externe, above=9mm of gui, minimum width=3.4cm] (gtk) {GTK~3, GLib\\[-1pt]
            \footnotesize GStreamer (optional)};

        \draw[fleche] (gui.south) -- (gui.south |- core.north);
        \draw[fleche] (cli.south) -- (core.north);
        \draw[fleche] (tests.south) -- (tests.south |- core.north);
        \draw[fleche, draw=ardoise, dashed] (gui.north) -- (gtk.south);

        \node[note, right=5mm of gtk, text width=6.4cm, align=left]
        {An arrow reads \enquote{depends on}.\\None leaves the engine.};
    \end{tikzpicture}
    \caption{The layers of the program. Three executables are built on the same engine.}
    \label{fig-couches}
\end{figure}

This rule is not a convention one promises to follow: it is \emph{enforced by the build}. The
engine is a static library whose only include path is the \fichier{include/} folder. The
continuous integration pipeline compiles and tests it in a container where GTK is not
installed (\cref{sec-ci}): if someone added a call to GTK to the engine, that step would stop
compiling.

Three programs rest on the engine. Having more than one is not a luxury. As long as there is
only one interface, nothing proves that the model is independent of it; the second interface
is the demonstration.

\section{From the course to the project: model, view, controller}

The \emph{model / view / controller} division was born at the end of the 1970s in the
Smalltalk environment~\cite{reenskaug1979, krasner1988}. It separates what a program
\emph{knows}, what it \emph{shows} and the way it \emph{reacts}.

The course describes how a graphical application works in three steps~\cite{toulouse2023}:
the user acts on the view, the corresponding action takes place in the model, then the view
is updated. It derives three libraries from this. \Cref{tab-mvc} shows how the project
implements them.

\begin{table}[ht]
    \centering
    \begin{tabularx}{\textwidth}{@{}l>{\raggedright\arraybackslash}X>{\raggedright\arraybackslash}X@{}}
        \toprule
                            & \textbf{Role (course)}                                                 & \textbf{Implementation (project)} \\
        \midrule
        \textbf{Model}      & User types and the functions that manipulate them: the game and its
        rules.              &
        \fichier{include/battleship/}, \fichier{src/core/} --- board, shots, deductions,
        turns, opponent.                                                                                                                 \\
        \addlinespace
        \textbf{View}       & Graphical components to be arranged in windows.                        &
        \fichier{src/ui/view.c} --- all the \emph{widgets}, no rule;
        \fichier{assets/style.css} for the appearance.                                                                                   \\
        \addlinespace
        \textbf{Controller} & Conductor: associates actions with events, synchronises the view.      &
        \fichier{src/ui/controller.c} --- asks the model to play, then resynchronises the
        view.                                                                                                                            \\
        \bottomrule
    \end{tabularx}
    \caption{Correspondence between the organisation of the course and that of the project.}
    \label{tab-mvc}
\end{table}

\subsection{Who owns what?}

The course discusses four ways of grouping the model and the view in the structure of the
controller, and keeps the one where the model is \emph{referenced} and the view \emph{owned}:
the model has an existence of its own, before and after any graphical interface, whereas the
view has no reason to exist without its controller. The project makes exactly this choice
(\cref{lst-controleur})\footnote{The listings of version~2 are abridged, and some of their
    comments are reworded for this report.}.

\begin{lstlisting}[caption={The structure of the controller (\fichier{src/ui/controller.c}).},
                   label={lst-controleur}]
struct bs_controller {
    GtkApplication *application;
    bs_game *model; /* borrowed: owned by the application */
    bs_view *view;  /* owned */
    /* ... */
};
\end{lstlisting}

The model, for its part, is declared by the application (\fichier{src/ui/app.c}), outside the
controller. It could be saved, replayed or handed to another interface without the graphical
interface knowing anything about it.

\subsection{The \enquote{function pointer} variant}
\label{sec-pointeurs}

In the course's model solution, the controller calls \code{vue\_tresor\_get\_boite()} to
fetch the buttons from the view and attaches the callbacks to them itself. A note adds that
another, more elegant way exists, at the price of function pointers. It is this variant that
the project implements.

The view connects its signals itself and translates each gesture of the user into an
\emph{intention}, expressed in the vocabulary of the game, which it passes on through a table
of callbacks (\cref{lst-actions}).

\begin{lstlisting}[caption={The intentions reported by the view (\fichier{src/ui/view.h}).},
                   label={lst-actions}]
typedef struct bs_view_actions {
    void (*level_chosen)(bs_difficulty level, gpointer user_data);
    void (*cell_activated)(bs_coord cell, gpointer user_data);
    void (*reshuffle_requested)(gpointer user_data);
    void (*restart_requested)(gpointer user_data);
    void (*quit_requested)(gpointer user_data);
    void (*game_over_answered)(bs_view_choice choice, gpointer user_data);
    void (*closed)(gpointer user_data);
} bs_view_actions;
\end{lstlisting}

The consequence is a strong one: \textbf{no \code{GtkWidget} crosses the interface of the
    view}. The controller does not handle buttons; it says \emph{what} to show --- \enquote{this
    cell is a hit} --- never \emph{how}. It could drive an entirely different graphical library
without a single line changing.

\section{The modules}

\begin{figure}[ht]
    \centering
    \begin{tikzpicture}[node distance=6mm]
        \node[boite] (types) {\textbf{types}\\[-1pt]\footnotesize coordinates,\\[-2pt]\footnotesize statuses};
        \node[boite, right=of types] (rng) {\textbf{rng}\\[-1pt]\footnotesize reproducible\\[-2pt]\footnotesize randomness};
        \node[boite, right=of rng] (board) {\textbf{board}\\[-1pt]\footnotesize fleet, shots,\\[-2pt]\footnotesize deductions};
        \node[boite, right=of board] (ai) {\textbf{ai}\\[-1pt]\footnotesize choice\\[-2pt]\footnotesize of a target};
        \node[boite, right=of ai] (game) {\textbf{game}\\[-1pt]\footnotesize turns,\\[-2pt]\footnotesize score, end};
        \draw[fleche] (rng) -- (types);
        \draw[fleche] (board) -- (rng);
        \draw[fleche] (ai) -- (board);
        \draw[fleche] (game) -- (ai);
    \end{tikzpicture}
    \caption{The modules of the engine form a chain: each one rests only on those before it.}
    \label{fig-modules}
\end{figure}

Inside the engine, the five modules form a strict chain (\cref{fig-modules}). Each follows
the pattern of the course: a header file that sets out the responsibilities --- data
structure, instructions for use, preconditions --- and a source file that implements them.
\Cref{tab-arborescence} gives the complete organisation of the repository.

\begin{table}[ht]
    \centering
    \small
    \begin{tabularx}{\textwidth}{@{}l>{\raggedright\arraybackslash}Xr@{}}
        \toprule
        \textbf{Folder}                                    & \textbf{Content}                                           & \textbf{Lines} \\
        \midrule
        \fichier{include/battleship/}, \fichier{src/core/} & Game engine: public interface and implementation           & \num{1343}     \\
        \fichier{src/cli/}                                 & Terminal interface and simulator                           & \num{645}      \\
        \fichier{src/ui/}                                  & GTK application: view, controller, shell, sound            & \num{1670}     \\
        \fichier{assets/}                                  & Images, sounds, style sheet (embedded in the executable)   & \num{199}      \\
        \fichier{tests/}                                   & Unit, property and terminal interface tests                & \num{1582}     \\
        \fichier{docs/}                                    & Architecture guide, decisions, this report                 &                \\
        \fichier{cmake/}, \fichier{packaging/}             & Build modules, desktop integration                         &                \\
        \fichier{.github/}                                 & Continuous integration and delivery                        &                \\
        \bottomrule
    \end{tabularx}
    \caption{Organisation of the repository. The lines count the code and its comments; for
        \fichier{assets/}, the style sheet alone.}
    \label{tab-arborescence}
\end{table}

\begin{retenir}
    The code of the tests (\num{1582} lines) is larger than that of the engine it verifies
    (\num{1343} lines, headers included). This ratio is usual for a component whose behaviour
    one wants to be able to guarantee.
\end{retenir}
